\subsection{矩阵的子式}

设 \(X\) 是 \((\xi_{ij})_{(i,j)\in I\times J}\) 型为 \((I,J)\) 的长方矩阵，其指标集 \(I\) 和 \(J\) 是全序的。若 \(H\subset I\) 与 \(K\subset J\) 是具有相同个数 \(p\) 个元素的有限子集，则存在唯一的递增双射 \(\phi :H\to K\) （集合论，III，§5，第 3 号，命题 6）；我们用 \(X_{H,K}\) 表示型为 \((H,H)\) 、等于 \((\xi_{i,\phi (j)})_{(i,j)\in H\times H}\) 的方阵。若 \(X\) 的元素属于交换环 \(A\) ，则行列式 \(\det (X_{H,K})\) 称为矩阵 \(X\) 的指标为 \(H,K\) 的子式；这些行列式（对所有有序对 \((H,K)\) 而言，其中 H 和 K 分别是 \(I\) 和 \(J\) 的含 \(p\) 个元素的子集）也称为 \(X\) 的 \(p\) 阶子式。利用这一记号：

命题9. 设 M 是一个具有基 \((e_i)_{i \in J}\) 的 A-模（有限或非有限），其指标集 J 是全序的。对每个整数 \(p > 0\) ，设 \((e_H)_{H \in \mathfrak{F}_p(J)}\) 为 \(\bigwedge^p(M)\) 的相应基（§7，第 8 号）。设 \((x_i)_{1 \leqslant i \leqslant p}\) 是 M 的 \(p\) 个元素组成的一个序列；设

\[
x _ {i} = \sum_ {j \in J} \xi_ {j i} e _ {j} \quad for i \in I = [1, p]
\]

并设 \(X\) 表示 \((\xi_{ji})\) 型为 (J, I) 的矩阵。则

\[
x _ {1} \wedge x _ {2} \wedge \dots \wedge x _ {p} = \sum_ {\mathrm{H} \in \mathfrak {F} _ {p} (J)} (\det X _ {\mathrm{H,I}}) e _ {\mathrm{H}}.\tag{17}
\]

其中 H 遍历集合 \(\mathfrak{F}_{p}(\mathrm{J})\) ，即 J 的具有 p 个元素的子集组成的集合。

这直接由第 4 号的公式 (12) 和第 3 号的公式 (6) 得出。

命题10. 设 M 和 N 是维数分别为 m 和 n 的两个自由 A-模， \(u: \mathbf{M} \to \mathbf{N}\) 是一个线性映射， \(X\) 是 \(u\) 关于 M 的基 \((e_i)_{1 \leqslant i \leqslant m}\) 与 N 的基 \((f_j)_{1 \leqslant j \leqslant n}\) 的矩阵。那么，对每个整数 \(p \leqslant \inf(m, n)\) ， \(\bigwedge^p(u)\) 关于基 \((e_{\mathbf{K}})_{\mathbf{K} \in \mathfrak{F}_p(\mathbf{I})}\) （ \(\bigwedge^p(\mathbf{M})\) 的基）与基 \((f_{\mathrm{H}})_{\mathrm{H} \in \mathfrak{F}_p(\mathrm{J})}\) （ \(\bigwedge^p(\mathbf{N})\) 的基）（其中我们记 \(\mathbf{I} = [1, m]\) 和 \(\mathbf{J} = [1, n]\)）的矩阵是矩阵 \((\det(X_{\mathrm{H},\mathbf{K}}))\) ，其型为 \((\mathfrak{F}_p(\mathbf{J}), \mathfrak{F}_p(\mathbf{I}))\) （因而有 \(\binom{n}{p}\) 行和 \(\binom{m}{p}\) 列）。

对 \(\mathbf{K} \subset \mathbf{J}\) 的含 \(p\) 个元素的子集，令 \((j_k)_{1 \leqslant k \leqslant p}\) 为 \(\mathbf{K}\) 的元素按递增顺序排列成的序列；根据 \(\bigwedge^p(u)\) 的定义及 §7 第 2 号公式 (4)

\[
\bigwedge^ {p} (u) \left(e _ {\mathbf {K}}\right) = u \left(e _ {j _ {1}}\right) \wedge u \left(e _ {j _ {2}}\right) \wedge \dots \wedge u \left(e _ {j _ {p}}\right).
\]

因此， \(\bigwedge^{p}(u)\) 的矩阵中位于指标H所在行与指标K所在列的元素，是元素 \(u(e_{j_1}) \wedge \cdots \wedge u(e_{j_p})\) 的指标H分量；因此由命题9，它等于 \(\det(X_{\mathrm{H},\mathbf{K}})\) 。

矩阵 \((\det(X_{\mathrm{H},\mathbf{K}}))\) 称为矩阵 X 的 p 次外幂，记作 \(\bigwedge^{p}(X)\) 。当 p = m = n 时， \(\bigwedge^{n}(X)\) 是仅含单个元素 \(\det(X)\) 的矩阵。

命题11. 设 M 是维数 n 为有限的自由 A-模；对 M 的每个自同态 u 以及 A 的每一对有序元素 \(\xi\) , \(\eta\) ，

\[
\det (\xi 1 _ {\mathrm{M}} + \eta u) = \sum_ {k \geqslant 0} \operatorname{Tr} \left(\bigwedge^ {k} (u)\right) \xi^ {n - k} \eta^ {k}.\tag{18}
\]

设 \((e_i)_{1\leqslant i\leqslant n}\) 是 M 的一组基，并令 \(\mathbf{I} = [1,n]\) ；为计算 (18) 的左边，我们必须构成乘积

\[
\left(\xi e _ {1} + \eta u (e _ {1})\right) \wedge \left(\xi e _ {2} + \eta u (e _ {2})\right) \wedge \dots \wedge \left(\xi e _ {n} + \eta u (e _ {n})\right)
\]

它等于诸项 \(\xi^{n - p}\eta^{p}z_{\mathbf{K}}\) 之和，其中

\[
z _ {K} = x _ {1} \wedge x _ {2} \wedge \dots \wedge x _ {n}
\]

其中 \(x_{i} = u(e_{i})\) （对 \(i \in \mathbf{K}\) ）， \(x_{i} = e_{i}\) （对 \(i \in \mathbf{H} = \mathbf{I} - \mathbf{K}\) ），这里整数 \(p\) 遍历区间 \([0, n]\) ，而对每个 \(p, \mathbf{K}\) 遍历 I 的具有 p 个元素的子集组成的集合。若 \(i_{1} < i_{2} < \cdots < i_{n-p}\) （相应地， \(j_{1} < j_{2} < \cdots < j_{p}\) ）是 H（相应地，K）的元素按递增顺序排列的结果，则我们可以写出（§7，第 8 号，定理 1 的推论 1 及公式 (19)）

\[
z _ {\mathbf {K}} = \rho_ {\mathrm{H}, \mathbf {K}} e _ {i _ {1}} \wedge e _ {i _ {2}} \wedge \dots \wedge e _ {i _ {n - p}} \wedge u (e _ {j _ {1}}) \wedge \dots \wedge u (e _ {j _ {p}}).
\]

但若 \(X\) 是 \(u\) 关于基 \((e_i)\) 的矩阵，则由命题 10

\[
u (e _ {j _ {1}}) \wedge \dots \wedge u (e _ {j _ {p}}) = \sum_ {\mathrm{L} \in \mathfrak {F} _ {p} (\mathrm{I})} (\det X _ {\mathrm{L}, \mathrm{K}}) e _ {\mathrm{L}}
\]

因此

\[
z _ {\mathbf {K}} = \rho_ {\mathrm{H}, \mathbf {K}} \sum_ {\mathbf {L} \in \mathfrak {F} _ {p} (\mathbf {I})} (\det X _ {\mathbf {L}, \mathbf {K}}) e _ {\mathrm{H}} \wedge e _ {\mathrm{L}}.
\]

现在，除 L = K 的情形外均有 \(H \cap L \neq \varnothing\) ；因此由 §7 第 8 号公式 (20) 得 \(z_{K} = (\det X_{K,K})e_{1} \wedge e_{2} \wedge \cdots \wedge e_{n}\) ，而公式 (18) 随后由命题 10 及矩阵的迹的定义（II，§10，第 11 号，公式 (49) 和 (50)）得出。

推论. 在与命题 11 相同的假设下，对自同态 \(\bigwedge(u)\) （A-模 \(\bigwedge(M)\) 的），

\[
\operatorname{Tr} (\bigwedge (u)) = \det (1 _ {\mathbf {M}} + u).\tag{19}
\]

只需在 (18) 中把 \(\xi\) 和 \(\eta\) 都换成 1，并注意到 \(\bigwedge(u)\) 关于由诸 \(e_{\mathbf{H}}\) （ \(\mathbf{H} \in \mathfrak{F}(\mathbf{I})\) ）组成的基的矩阵，是由诸 \(\bigwedge^{k}(u)\) （ \(k \geqslant 0\) ）的矩阵组成的分块对角矩阵（II，§ 10，第 7 号，例 IV）。

